\begin{Exercise}[label={ex:ricerca}]
    Si realizzi un programma che legga da input un numero non noto a priori di interi (al più 10) terminati da 0.
    Tale sequenza potrà evantualmente contenere numeri ripetuti.
    A tal fine, si realizzi una funziona \texttt{leggi(int v[], int dim)} che riceva in ingresso un vettore (vuoto) e la sua dimensione fisica,
    e restituisca la dimensione logica.
    
    Si chieda poi l'inserimento di un valore da cercare nell'array. 
    Si realizzi una funzione \texttt{search(int v[], int dim, int k)} che stampi a video tutti gli indici relativi alle posizioni in cui il numero k sia presente nell’array. 
\end{Exercise}

\begin{Answer}[ref={ex:ricerca}]
  	\lstinputlisting[style=cstylesmall]{src/ricerca.c}
\end{Answer}